\subsection{Desarrollo}
\label{subsec:desarrollo}

\subsubsection{Materiales e Instrumentos}

A continuación se detallan los materiales utilizados en cada práctica:

\todo{Completar la incertidumbre de cada instrumento según sea necesario.}

\begin{table}[H]
    \centering
    \label{tab:instrumentos}
    \begin{tabular}{|p{5cm}|p{4cm}|p{6.5cm}|}
        \hline
        \textbf{Instrumento} & \textbf{Incertidumbre} & \textbf{Uso en la práctica} \\ \hline
        \csvreader[
            separator=semicolon,
            head=true,
            late after line=\\ \hline
        ]{secciones/desarrollo/instrumentos.csv}{}
        {\csvcoli & \csvcolii & \csvcoliii}
    \end{tabular}
\end{table}

\subsubsection{Procedimiento}

\paragraph{Práctica 1 --- Ley de Ohm}

\todo{Completar con el bosquejo del procedimiento.}

\paragraph{Práctica 2 --- Medición de resistencias con multímetro}
Se midieron individualmente tres resistencias de la caja de laboratorio utilizando el multímetro digital UT58D en modo ohmímetro.
Para cada resistencia se identificaron las franjas de color para determinar el valor nominal y la tolerancia.
Se seleccionó el rango del multímetro adecuado (el menor rango que contuviera el valor esperado) y se conectaron las puntas a los extremos de cada resistencia.
Para R3 se realizó primero una lectura en rango 2000 k$\Omega$ y luego se ajustó al rango 20 k$\Omega$ para obtener mayor precisión.
Se registró el valor indicado en el display y se calculó la incerteza $\Delta R$ según la Ecuación~\ref{eq:incerteza-multimetro} con los parámetros del rango utilizado.

\paragraph{Práctica 3 --- Influencia del aparato de medida}

\todo{Completar con el bosquejo del procedimiento.}
